\begin{table}[htbp]
\raggedright
\caption{Profile moderators of attacked effectivity: controlled contrasts, SEM paths, and descriptive susceptibility weights.}
\label{tab:multivariate}
{
\footnotesize
\setlength{\tabcolsep}{4pt}
\renewcommand{\arraystretch}{1.08}
\resizebox{\linewidth}{!}{%
\begin{tabular}{p{0.34\linewidth}rrrrrrrrr}
\toprule
Moderator & Role & Controlled mean b & Controlled p & Ridge mean b & Weight \% & SEM mean b & SEM mean |b| & SEM min p & SEM paths \\
\midrule
Big Five Conscientiousness Mean \% & core & -1.034 & 0.123 &  &  & -1.185 & 2.010 & 0.061 & 8.000 \\
Big Five Neuroticism Mean \% & core & -0.934 & 0.134 &  &  & -1.039 & 1.560 & 0.095 & 8.000 \\
Sex Other & core & -0.429 & 0.780 & 0.002 & 0.779 & -0.776 & 1.952 & 0.237 & 8.000 \\
Big Five Openness To Experience Mean \% & core & 0.077 & 0.916 &  &  & -0.016 & 1.391 & 0.114 & 8.000 \\
Sex Female & core & 0.485 & 0.741 & -0.000 & 0.534 & -0.259 & 2.512 & 0.162 & 8.000 \\
Big Five Extraversion Mean \% & exploratory & 1.757 & 0.012 &  &  &  &  &  &  \\
Big Five Agreeableness Mean \% & exploratory & -0.078 & 0.910 &  &  &  &  &  &  \\
\bottomrule
\end{tabular}
}
}
\par\smallskip
{\normalsize Note. Controlled coefficients summarize moderator associations with mean attacked effectivity. SEM columns summarize the repeated-outcome path model across the attacked opinion leaves. Weight percentages come from the target-conditional regularized aggregation used to compute the post hoc susceptibility index.}
\end{table}